\documentclass[11pt,letterpaper]{article}
\usepackage[T1]{fontenc}
\usepackage{lmodern,microtype}
\usepackage[margin=0.94in,headheight=14pt]{geometry}
\usepackage{amsmath,amssymb,amsthm,booktabs,array,fancyhdr}
\usepackage{xcolor}
\usepackage[colorlinks=true,linkcolor=black,citecolor=black,urlcolor=blue!55!black]{hyperref}
\hypersetup{pdftitle={Exact modular periods of Fibonacci-weighted Stirling transforms},
  pdfauthor={Blake Stone, with OpenAI Codex}}
\pagestyle{fancy}\fancyhf{}
\lhead{\small Exact modular periods}\rhead{\small Fibonacci-weighted Stirling transforms}
\cfoot{\thepage}\renewcommand{\headrulewidth}{0.3pt}
\newtheorem{theorem}{Theorem}[section]
\newtheorem{lemma}[theorem]{Lemma}
\newtheorem{corollary}[theorem]{Corollary}
\newcommand{\FF}{\mathbb F}
\newcommand{\ZZ}{\mathbb Z}
\newcommand{\lcm}{\operatorname{lcm}}
\newcommand{\Li}{\operatorname{Li}}
\newcommand{\ord}{\operatorname{ord}}
\newcommand{\stir}[2]{\genfrac{\{}{\}}{0pt}{}{#1}{#2}}
\newcommand{\coef}[1]{[X^{#1}]}
\setlength{\parskip}{0.22em}
\setlength{\abovedisplayskip}{7pt plus 2pt minus 1pt}
\setlength{\belowdisplayskip}{7pt plus 2pt minus 1pt}
\raggedbottom

\begin{document}
\thispagestyle{plain}
\begin{center}
{\LARGE\bfseries Exact modular periods of\\[2pt] Fibonacci-weighted Stirling transforms\par}
\vspace{0.6em}
{\large Blake Stone\par}
{\small Research manuscript prepared with OpenAI Codex\par}
\end{center}

\begin{abstract}
Kinlaw, Morris and Thiagarajan conjectured that the least eventual period of
the first of two sequences attached to the golden-ratio polylogarithm equals
Carmichael's function at every modulus. They also conjectured a family of
half-periods for the second sequence. We prove both assertions and classify
the exceptional moduli of the second sequence exactly. Our proof isolates a
nonzero coefficient in the difference between moments separated by $p-1$,
then lifts this obstruction through all odd prime powers. We also determine
the least eventual periods of two central-binomial Stirling transforms at
every modulus.
\end{abstract}

\section{The period conjectures}

Let $\phi=(1+\sqrt5)/2$. For $n\ge1$, the two integers $A_n,B_n$ are defined
by $\Li_{-n}(1/\phi)=A_n+B_n\phi$. They appear in Ledin's Fibonacci power
sums and in OEIS A000556 and A000557 \cite{KMT,OEIS}. Proposition~4.10 of
Kinlaw, Morris and Thiagarajan gives
\begin{equation}\label{eq:stirling}
 A_n=\sum_{j=0}^n \stir nj j!F_{j+1},\qquad
 B_n=\sum_{j=0}^n \stir nj j!F_{j+2}.
\end{equation}
Here $F_0=0,F_1=1$ and the braces denote Stirling numbers of the second kind.
We extend both sequences by $A_0=B_0=1$ using \eqref{eq:stirling}. This
convention preserves every eventual period.

For an integer sequence $S$, let $P_S(m)$ be its least positive eventual
period modulo $m$. Write $\lambda(m)$ for the exponent of the group of units
modulo $m$, with $\lambda(1)=1$. Theorem~4.14 of \cite{KMT} establishes
$P_A(m),P_B(m)\mid\lambda(m)$. Remark~4.15 asks whether the first period
always attains this bound and proposes specific infinite families where the
second is half the bound.

\begin{theorem}\label{th:main}
For every positive modulus $m$, $P_A(m)=\lambda(m)$. Write $m=2^e q$ with
$q$ odd. Then $P_B(m)=\lambda(m)/2$ exactly when either condition below holds.
\[
 m\in\{3,6\},\qquad\text{or}\qquad
 e\ge4\ \text{ and }\ v_2(\lambda(q))<e-2.
\]
At every other modulus $P_B(m)=\lambda(m)$.
\end{theorem}

For $e\ge4$ this includes $m=2^e$ and $m=3\cdot2^e$ as conjectured.
It accounts for $112,144,160,176$ in the authors' numerical list and adds
$3,6$. The latter two have a short direct check. For $n\ge1$ the
$j=1$ summand for $B_n$ is $2$, the $j=2$ summand is divisible by $6$, and
$6\mid j!$ for $j\ge3$. Thus $B_n\equiv2\pmod6$ throughout.

\clearpage
\section{Finite power sums and lifting at odd primes}

Suppose $S_n=n![t^n]G(e^t-1)$ with $G(z)=\sum_{j\ge0}g_jz^j\in\ZZ[[z]]$.
For $M\mid J!$ set
\[
 K_J(X)=\sum_{j=0}^{J-1}g_j(X-1)^j=\sum_{u=0}^{J-1}c_uX^u.
\]
The expansion into Stirling numbers and then the binomial theorem give
\begin{equation}\label{eq:spectrum}
 S_n\equiv\sum_{u=0}^{J-1}c_u u^n\pmod M .
\end{equation}
Indeed the exponential coefficient of $(e^t-1)^j$ is $j!\stir nj$, a
multiple of $M$ for $j\ge J$. For $M=p^e$ the nonunit bases give zero at
$n\ge e$, and all unit powers have period dividing $\lambda(p^e)$. The least
period divides every period, by the division algorithm on a periodic tail.
The Chinese remainder theorem gives the exact reduction
\begin{equation}\label{eq:crt}
 P_S(m)=\lcm_{p^e\parallel m}P_S(p^e).
\end{equation}
These upper bounds have precedents in the Fibonacci-weighted case \cite{KMT},
the ordered Bell case \cite{Poonen}, and the general integral-substitution
setting \cite{Danesh}. We use \eqref{eq:spectrum} also to obtain lower bounds.

\begin{lemma}[Odd primes]\label{lem:odd}
Let $p$ be odd. In $\FF_p[X]$ write, uniquely,
\[
 K_{2p}(X)=U(X)+(X^p-1)V(X),\qquad \deg U,\deg V<p .
\]
Put $u_r=\coef r U$ and $v_r=\coef r V$.
\begin{enumerate}
\item The least period modulo $p$ is the exponent of the subgroup of
$\FF_p^\times$ generated by $\{r:1\le r<p,\ u_r\ne0\}$.
\item If $v_1\ne0$, shift $p-1$ fails modulo $p^2$ on the periodic tail.
\item Whenever shift $p-1$ fails modulo $p^2$, one has
$p^{e-1}\mid P_S(p^e)$ for every $e\ge2$.
\end{enumerate}
\end{lemma}

\begin{proof}
Modulo $p$ the terms in \eqref{eq:spectrum} with bases $r$ and $r+p$
combine into $u_rr^n$. On any $p-1$ consecutive indices the columns
$r^n$, $1\le r<p$, are independent over $\FF_p$ by the Vandermonde
determinant, up to multiplication of columns by units. Thus a tail shift
$L$ works exactly when $r^L=1$ at each nonzero coefficient.

Use $J=2p$ also modulo $p^2$, since $p^2\mid(2p)!$. At $n\ge2$ divide
$S_{n+p-1}-S_n$ by $p$ and reduce modulo $p$. Each unit base $u$
contributes its coefficient times $u^n(u^{p-1}-1)/p$. Grouping by residue
again gives a sum of independent columns. In the $r=1$ column the base
$1$ contributes zero. The base $1+p$ has coefficient $v_1$ and
$((1+p)^{p-1}-1)/p\equiv-1\pmod p$. The combined coefficient is $-v_1$,
which proves the second assertion.

For $j\ge0$, a unit $u$ satisfies
$u^{(p-1)p^j}-1\equiv p^j(u^{p-1}-1)\pmod{p^{j+2}}$.
To see this, write $u^{p-1}=1+pb$ and successively take $p$th powers.
For odd $p$ the binomial expansion gives $(1+p^a b)^p\equiv1+p^{a+1}b\pmod{p^{a+2}}$
for $a\ge1$. Apply \eqref{eq:spectrum} at the common modulus $p^{j+2}$.
Nonunit terms vanish for $n\ge j+2$, so
\begin{equation}\label{eq:oddlift}
 S_{n+(p-1)p^j}-S_n\equiv p^j(S_{n+p-1}-S_n)\pmod{p^{j+2}}.
\end{equation}
Failure of shift $p-1$ modulo $p^2$ gives arbitrarily large indices of
failure, hence so does shift $(p-1)p^{e-2}$ modulo $p^e$. Every divisor of
$\lambda(p^e)$ missing a factor of $p$ from $p^{e-1}$ divides that shift.
\end{proof}

\clearpage
\section{Fibonacci weights at primes \texorpdfstring{$p\ge7$}{p at least 7}}

The two functions $R(X)=G(X-1)$ coming from \eqref{eq:stirling} are
\[
 R_A(X)=\frac{1}{Q(X)},\qquad
 R_B(X)=\frac{X}{Q(X)},\qquad Q(X)=1+X-X^2.
\]
In characteristic $p$, $(X-1)^p=X^p-1$. Thus $U$ in Lemma~\ref{lem:odd}
is the representative of $R$ of degree less than $p$ modulo $X^p-1$.
The denominator is invertible there since $Q(1)=1$. Write $Q U\equiv N$,
where $N=1$ or $X$. In degrees $3\le r\le p-1$ this reads
\[
 u_r=u_{r-2}-u_{r-1}.
\]
If two of $u_1,\ldots,u_{p-1}$ vanished at distance one or two, the recurrence
would give adjacent zeros and then propagate zeros in both directions
throughout that list. This would leave $U$ constant, which is impossible
since no constant times $Q$ equals $N$ modulo $X^p-1$. Thus at most
$\lceil(p-1)/3\rceil$ of these $p-1$ coefficients vanish. For $p\ge7$ the
nonzero support has more than $(p-1)/2$ elements, so it generates the entire
multiplicative group. Lemma~\ref{lem:odd} gives period $p-1$ modulo $p$.

To compute $v_1$, factor $Q(X)=(1-aX)(1-bX)$ over $\FF_{p^2}$. The
distinct elements $a,b$ have $a+b=-1$ and $ab=-1$. Put $\delta=1$ for $A$
and $\delta=0$ for $B$. Partial fractions give
\[
 R(X)=\frac{1}{a-b}\left(\frac{a^\delta}{1-aX}
                          -\frac{b^\delta}{1-bX}\right).
\]
Introduce $Y=X^p$. Using
$(1-aX)^{-1}=(1+aX+\cdots+a^{p-1}X^{p-1})/(1-a^pY)$,
the coefficient of $X^r$ in the basis $1,X,\ldots,X^{p-1}$ over rational
functions in $Y$ is
\[
 W_r(Y)=\frac{1}{a-b}\left(\frac{a^{r+\delta}}{1-a^pY}
                       -\frac{b^{r+\delta}}{1-b^pY}\right).
\]
The denominators are nonzero at $Y=1$. Reducing modulo
$(Y-1)^2=(X-1)^{2p}$ identifies $W_r(1)$ with $u_r$ and $W'_r(1)$
with $v_r$. Since each of $a^p,b^p$ satisfies $1-z=z^2$,
\[
 v_1=\frac{a^{1+\delta}(a^p)^{-3}-b^{1+\delta}(b^p)^{-3}}{a-b}.
\]
Frobenius either fixes both roots or interchanges them. In the first case
this is $(a^{\delta-2}-b^{\delta-2})/(a-b)=1$ for both values of $\delta$.
In the second case $b=-1/a$ gives
$v_1=-(a^{4+\delta}-b^{4+\delta})/(a-b)$.
The quotients $(a^j-b^j)/(a-b)$ at $j=1,\ldots,5$ are
$1,-1,2,-3,5$ by the quadratic recurrence. Thus the possible values are
\[
 \begin{array}{c|cc}
 & A & B\\ \hline
 \text{roots fixed} & 1 & 1\\
 \text{roots interchanged} & -5 & 3
 \end{array}
\]
and all are nonzero for $p\ge7$. Lemma~\ref{lem:odd} now forces both
periods to equal $(p-1)p^{e-1}$ at every $p^e$ with $p\ge7$.

\clearpage
\section{Small primes and completion of the classification}

At the primes $3$ and $5$ the following are the least periods, obtained
from the finite sum \eqref{eq:spectrum}. Section~\ref{sec:cert} gives the
exact finite test and verification details.
\[
\begin{array}{c|rrrr}
 M&3&9&5&25\\\hline
 P_A(M)&2&6&4&20\\
 P_B(M)&1&6&4&20
\end{array}
\]
At both prime squares the period of each sequence contains the odd prime and
the entire factor $p-1$. Lemma~\ref{lem:odd}(3) gives the full $p$-part at
every higher power, and reduction to the square preserves the lower bound
$p-1$ on the prime-to-$p$ part.

\begin{lemma}[Powers of two]\label{lem:two}
For $j\ge1$ let $w_j$ be the largest $t$ such that $2^j$ is an eventual
period modulo $2^t$, when this largest $t$ exists. If $w_j+1<2j+4$, then
$w_{j+1}=w_j+1$. The same test continues to apply at every later $j$.
\end{lemma}

\begin{proof}
Write $\Delta_h S_n=S_{n+h}-S_n$. For an odd base $u$,
$2^{j+2}\mid u^{2^j}-1$ by $u^2\equiv1\pmod8$ and repeated squaring.
The difference $\Delta_{2h}u^n-2\Delta_hu^n$ equals
$u^n(u^h-1)^2$. With $h=2^j$, \eqref{eq:spectrum} therefore gives
$\Delta_{2h}S_n\equiv2\Delta_hS_n\pmod{2^{2j+4}}$ on the tail.
The right side vanishes modulo $2^{w_j+1}$ and fails to vanish modulo
$2^{w_j+2}$ arbitrarily far along the tail. The congruence transfers both
facts. The bound $2j+4$ grows by two at each step while $w_j$ grows by one.
\end{proof}

At modulus $32$, a full tail block starting at $n=5$ is
\[
\begin{array}{c|rrrrrrrr}
 n&5&6&7&8&9&10&11&12\\\hline
 A_n\bmod32&15&25&23&1&31&9&7&17\\
 B_n\bmod32& 2& 0&18&0& 2&0&18& 0
\end{array}
\]
This is a full block because $\lambda(32)=8$. Reading differences two
steps apart gives $w_1=3$ for $A$ and $w_1=4$ for $B$. Lemma~\ref{lem:two}
therefore gives $w_j=j+2$ for $A$ and $w_j=j+3$ for $B$, for every $j\ge1$.
Reduction of the same block handles the initial powers of two. Since each
period at a power of two is itself a power of two, these thresholds give
\[
 P_A(2^e)=\lambda(2^e),\qquad
 P_B(2^e)=
 \begin{cases}
 \lambda(2^e),&1\le e\le3,\\
 2^{e-3},&e\ge4 .
 \end{cases}
\]

\begin{proof}[Completion of the proof of Theorem~\ref{th:main}]
Equation~\eqref{eq:crt} gives the claim for $A$. For odd $q$, the only
local loss for $B$ is at $3^1$, where the missing factor is two. Every
other odd prime-power period is even. Thus $P_B(q)=\lambda(q)$ except when
$q=3$. At the moduli $3$ and $6$ the period is one. If $e=2,3$ the factor
two from $2^e$ restores the full period even when $q=3$. For $e\ge4$ the
power of two contributed by $2^e$ drops from $2^{e-2}$ to $2^{e-3}$.
The least common multiple drops by a factor two exactly when the odd
part's period supplies fewer than $e-2$ factors of two. This is equivalent
to $v_2(\lambda(q))<e-2$, including $q=1$ and $q=3$.
\end{proof}

\clearpage
\section{Two central-binomial transforms}

For $a=1,2$ define
\[
 S_n(a)=\sum_{j=0}^n\stir nj j!\binom{2j}{j}a^j,\qquad
 G_a(z)=(1-4az)^{-1/2}.
\]
Write $C=S(1)$ and $D=S(2)$, OEIS A354242 and A354253. Bala's comments
in these entries ask for an eventual period dividing Euler's totient
\cite{OEIS,Danesh}. The same lifting method determines the least period for
each modulus, including a change at $3^6$ for $D$.

\begin{theorem}\label{th:central}
For both sequences the least period at $p^e$ is $(p-1)p^{e-1}$ for every
prime $p\ge5$. At $p=3$ one has
\[
\begin{aligned}
P_C(3^e)&=\begin{cases}1,&e=1,\\2\cdot3^{e-1},&e\ge2,\end{cases}
&
P_{D}(3^e)&=\begin{cases}3^{e-1},&1\le e\le5,\\2\cdot3^{e-1},&e\ge6.\end{cases}
\end{aligned}
\]
At powers of two the periods are
\[
\begin{aligned}
P_C(2^e)&=\begin{cases}1,&e\le3,\\2,&4\le e\le6,\\2^{e-5},&e\ge7,\end{cases}
&
P_{D}(2^e)&=\begin{cases}1,&e\le6,\\2,&7\le e\le8,\\2^{e-7},&e\ge9.\end{cases}
\end{aligned}
\]
For arbitrary $m$ take the least common multiple over its prime powers.
\end{theorem}

\begin{proof}
Let $p\ge7$ and $h=(p-1)/2$. In $\FF_p[[z]]$ we have
$G_a(z)=(1-4az)^hG_a(z^p)$ by Frobenius and the defining square.
Through degrees below $2p$ this is $(1-4az)^h(1+2az^p)$. With $z=X-1$ this
yields the following polynomials in Lemma~\ref{lem:odd}.
\[
 U(X)=(1+4a-4aX)^h,\qquad V(X)=2aU(X).
\]
For $a=1,2$ all coefficients of $U$ through degree $h$ are nonzero, since
$p$ divides none of $a,1+4a,h!$. The subgroup generated by $1,\ldots,h$
contains $2h=-1$ in $\FF_p$, as it contains both $2$ and $h$. It therefore
contains $\pm1,\ldots,\pm h$, the whole unit group. Moreover $v_1=2au_1\ne0$.
Lemma~\ref{lem:odd} gives the full period at all such prime powers.

The finite tests of Section~\ref{sec:cert} give period $4$ modulo $5$
and $20$ modulo $25$ for both sequences. At powers of three the initial
periods for $C$ are $1,6$. For $D$ the first six are
$1,3,9,27,81,486$. The shift $2$ fails modulo $9$ for both sequences.
Lemma~\ref{lem:odd}(3) forces the full power of three for all $e\ge2$.
Reduction to $9$ supplies the factor two for $C$, and reduction to $729$
supplies it for $D$ from $e=6$ onward. The argument for $5$ uses the
two periods at $5,25$ in the same way.

At the prime two, the initial periods for $C$ through $2^9$ and for
$D$ through $2^{13}$ are
\[
\begin{array}{c|rrrrrrrrrrrrr}
e&1&2&3&4&5&6&7&8&9&10&11&12&13\\\hline
C&1&1&1&2&2&2&4&8&16&&&&\\
D&1&1&1&1&1&1&2&2&4&8&16&32&64
\end{array}
\]
It follows that $w_3(C)=8$ and $w_5(D)=12$. Both satisfy the strict
inequality of Lemma~\ref{lem:two}, so all subsequent $w_j$ increase by one
per step. This gives precisely the stated powers of two. Equation~\eqref{eq:crt}
handles arbitrary moduli.
\end{proof}

\clearpage
\section{Reproducing the finite checks}\label{sec:cert}

All finite moduli used in the proofs divide $18!$. Here is a direct
integer certificate for each table entry. Form the coefficients of
$K_{18}(X)$ in \eqref{eq:spectrum} by Horner's rule, using successively
$P\gets g_j+(X-1)P$ for $j=17,16,\ldots,0$, starting with $P=0$.
For a candidate period $T$ and $M=p^e$, form
\[
 b_u=c_u u^e(u^T-1)\pmod M,\qquad
 d_i=\sum_{u=0}^{17} b_u u^i\pmod M
 \quad(0\le i<18).
\]
Then $T$ is a period exactly when all eighteen $d_i$ vanish. Indeed these
are the first eighteen differences starting at $n=e$. The infinite
difference sequence satisfies the monic recurrence with characteristic
polynomial $\prod_{u=0}^{17}(Z-u)$, so eighteen consecutive zeros force
every later difference to vanish. Conversely \eqref{eq:spectrum} gives
periodicity starting at $e$, so any eventual period holds there already.
To certify minimality, repeat the test at $T/\ell$ for each prime divisor
$\ell$ of $T$ and check that at least one residue is nonzero for each test.
All candidates in the finite tables have prime divisors among $2,3,5$.

For example, modulo $9$ the tail block of $B$ from $n=2$ is
$8,5,2,2,5,8$. Shifts $2$ and $3$ both fail at its first entry and
$6=\lambda(9)$ is a period. The companion file \texttt{SmallCases.lean}
checks all finite inputs to the two theorems by kernel reduction using
the preceding test. The program \texttt{verify.py} also computes complete
blocks of length $\lambda(p^e)$ at the finite base cases. For an independent
check it differentiates the defining generating functions and tests the
resulting terms at every modulus up to $500$.

\begin{thebibliography}{9}
\bibitem{KMT} P. Kinlaw, M. Morris and S. Thiagarajan.
\textit{Sums related to the Fibonacci sequence}.
Fibonacci Quarterly \textbf{60} (2022), 136--150.
\href{https://www.fq.math.ca/Papers/60-2/kinlaw03042021.pdf}{Publisher's PDF}.
See Proposition 4.10, Theorem 4.14 and Remark 4.15.

\bibitem{Poonen} B. Poonen. \textit{Periodicity of a combinatorial sequence}.
Fibonacci Quarterly \textbf{26} (1988), 70--76.
\href{https://www.mathstat.dal.ca/FQ/Scanned/26-1/poonen.pdf}{Publisher's archive}.

\bibitem{Danesh} P. Danesh. \textit{Eventual periodicity of weighted
Stirling transforms and two conjectures of Bala}. Cambridge Open Engage,
2026, preprint. \href{https://doi.org/10.33774/coe-2026-4d6jr}{doi:10.33774/coe-2026-4d6jr}.

\bibitem{OEIS} OEIS Foundation. The On-Line Encyclopedia of Integer Sequences.
\href{https://oeis.org/A000556}{A000556},
\href{https://oeis.org/A000557}{A000557},
\href{https://oeis.org/A354242}{A354242} and
\href{https://oeis.org/A354253}{A354253}.
For the central-binomial questions see P. Bala's comments dated July 2022.
\end{thebibliography}
\end{document}
